\documentclass[10pt,a4paper]{amsart}
\usepackage{amsmath,amssymb,amsthm,graphicx}
% Retained preamble for the Stacks Project chapter snapshot in this directory.
% The chapter's own preamble.tex is a copy of the shared Stacks Project preamble; the block
% below is sliced MECHANICALLY, line for line, from it. Nothing was added, removed or
% redefined; the wrapper class, the xr-hyper/hyperref wiring and the externaldocument list
% are omitted because the assembler supplies the document class.
%
% lines taken from preamble.tex: 3, 5, 9, 11, 12, 156, 157, 158, ... (51 lines total)
%
% package coverage restored from the chapter's own preamble (2026-10-13): xy, xy-2cell

\usepackage[all]{xy}
\xyoption{2cell}
\UseAllTwocells

\newenvironment{reference}{\comment}{\endcomment}
\newenvironment{slogan}{\comment}{\endcomment}
\newenvironment{history}{\comment}{\endcomment}
\theoremstyle{plain}
\newtheorem{theorem}[subsection]{Theorem}
\newtheorem{proposition}[subsection]{Proposition}
\newtheorem{lemma}[subsection]{Lemma}
\theoremstyle{definition}
\newtheorem{definition}[subsection]{Definition}
\newtheorem{example}[subsection]{Example}
\newtheorem{exercise}[subsection]{Exercise}
\newtheorem{situation}[subsection]{Situation}
\theoremstyle{remark}
\newtheorem{remark}[subsection]{Remark}
\newtheorem{remarks}[subsection]{Remarks}
\numberwithin{equation}{subsection}
\def\lim{\mathop{\mathrm{lim}}\nolimits}
\def\colim{\mathop{\mathrm{colim}}\nolimits}
\def\Spec{\mathop{\mathrm{Spec}}}
\def\Hom{\mathop{\mathrm{Hom}}\nolimits}
\def\Ext{\mathop{\mathrm{Ext}}\nolimits}
\def\SheafHom{\mathop{\mathcal{H}\!\mathit{om}}\nolimits}
\def\SheafExt{\mathop{\mathcal{E}\!\mathit{xt}}\nolimits}
\def\Sch{\mathit{Sch}}
\def\Mor{\mathop{\mathrm{Mor}}\nolimits}
\def\Ob{\mathop{\mathrm{Ob}}\nolimits}
\def\Sh{\mathop{\mathit{Sh}}\nolimits}
\def\NL{\mathop{N\!L}\nolimits}
\def\CH{\mathop{\mathrm{CH}}\nolimits}
\def\proetale{{pro\text{-}\acute{e}tale}}
\def\etale{{\acute{e}tale}}
\def\QCoh{\mathit{QCoh}}
\def\Ker{\mathop{\mathrm{Ker}}}
\def\Im{\mathop{\mathrm{Im}}}
\def\Coker{\mathop{\mathrm{Coker}}}
\def\Coim{\mathop{\mathrm{Coim}}}
\DeclareMathSymbol{\boxtimes}{\mathbin}{AMSa}{"02}
\def\QCohstack{\mathcal{QC}\!\mathit{oh}}
\def\Cohstack{\mathcal{C}\!\mathit{oh}}
\def\Spacesstack{\mathcal{S}\!\mathit{paces}}
\def\Quotfunctor{\mathrm{Quot}}
\def\Hilbfunctor{\mathrm{Hilb}}
\def\Curvesstack{\mathcal{C}\!\mathit{urves}}
\def\Polarizedstack{\mathcal{P}\!\mathit{olarized}}
\def\Complexesstack{\mathcal{C}\!\mathit{omplexes}}
\def\Pic{\mathop{\mathrm{Pic}}\nolimits}
\def\Picardstack{\mathcal{P}\!\mathit{ic}}
\def\Picardfunctor{\mathrm{Pic}}
\def\Deformationcategory{\mathcal{D}\!\mathit{ef}}
\usepackage{cleveref}
\crefname{originaldefinition}{original definition}{original definitions}
\Crefname{originaldefinition}{Original definition}{Original definitions}
\makeatletter
\newlabel{groupoids-definition-equivalence-relation}{{022P}{source}{}{sourcepointer-1}{}}
\newlabel{groupoids-definition-equivalence-relation@cref}{{[originaldefinition][0][]022P}{[1][0][]source}{}{}{}}
\makeatother
\crefname{originaldefinition}{original definition}{original definitions}
\Crefname{originaldefinition}{Original definition}{Original definitions}
\makeatletter
\newlabel{sites-definition-sheaves-injective-surjective}{{00WM}{source}{}{sourcepointer-2}{}}
\newlabel{sites-definition-sheaves-injective-surjective@cref}{{[originaldefinition][0][]00WM}{[1][0][]source}{}{}{}}
\makeatother
\crefname{originallemma}{original lemma}{original lemmas}
\Crefname{originallemma}{Original lemma}{Original lemmas}
\makeatletter
\newlabel{sites-lemma-coequalizer-surjection}{{086K}{source}{}{sourcepointer-3}{}}
\newlabel{sites-lemma-coequalizer-surjection@cref}{{[originallemma][0][]086K}{[1][0][]source}{}{}{}}
\makeatother
\title{A surjective etale chart gives an etale equivalence relation with a coequalizer diagram}
\author{The Stacks Project Authors}
\date{stacks-project:spaces.tex:local-snapshot-20261012; GNU FDL 1.2 or later}
\begin{document}
\maketitle
\setcounter{subsection}{1}
\noindent\textit{Editorial note.} All delivered mathematical text is the author's own, copied from the pinned Stacks Project chapter file spaces.tex; no statement or proof step was rewritten, reordered or removed.\par\medskip
\noindent\textit{Editorial note.} The statement contains an xymatrix display; it is part of the counted statement span and is preserved verbatim, with the chapter's own Xy-pic setup retained in the preamble.\par\medskip
\noindent\textit{Editorial note.} The three references of the counted spans are recorded as non-load-bearing navigation pointers; the argument that consumes them stays inside the counted proof span.\par\medskip
\noindent\textit{Editorial note.} Preamble: sources/stacks-spaces/retained-preamble.tex, sliced mechanically from the chapter's own preamble.tex.\par\medskip
\noindent\textit{Allowed standard prerequisites (editorial).}\begin{itemize}
\item Definition of an equivalence relation on a category. Groupoids, Definition groupoids-definition-equivalence-relation, used at proof lines 1393-1395. Supplies the relation axioms checked on T-valued points.
\item Definition of injective and surjective maps of sheaves. Sites, Definition sites-definition-sheaves-injective-surjective, used at proof lines 1417-1419. Is the sense in which U -> F is shown surjective.
\item Surjectivity implies the coequalizer statement. Sites, Lemma sites-lemma-coequalizer-surjection, used at proof lines 1424-1426. Concludes the coequalizer diagram.
\end{itemize}\medskip
\ifdefined\hypertarget\hypertarget{sourcepointer-1}{}\fi
\noindent\textit{Source reference.} Original definition 022P: focused-seeds/AGP/sources/stacks-snapshot-20261011/tags.txt, source lines 2714--2714. Navigation pointer into the Stacks Project; the standard fact it supplies is applied inside the counted proof span and the argument resting on it stays there.\par\medskip
\ifdefined\hypertarget\hypertarget{sourcepointer-2}{}\fi
\noindent\textit{Source reference.} Original definition 00WM: focused-seeds/AGP/sources/stacks-snapshot-20261011/tags.txt, source lines 1188--1188. Navigation pointer into the Stacks Project; the standard fact it supplies is applied inside the counted proof span and the argument resting on it stays there.\par\medskip
\ifdefined\hypertarget\hypertarget{sourcepointer-3}{}\fi
\noindent\textit{Source reference.} Original lemma 086K: focused-seeds/AGP/sources/stacks-snapshot-20261011/tags.txt, source lines 10395--10395. Navigation pointer into the Stacks Project; the standard fact it supplies is applied inside the counted proof span and the argument resting on it stays there.\par\medskip
\begin{lemma}
\label{lemma-space-presentation}
Let $F$ be an algebraic space over $S$. Let $f : U \to F$ be a
surjective \'etale morphism from a scheme to $F$. Set $R = U \times_F U$.
Then
\begin{enumerate}
\item $j : R \to U \times_S U$ defines an equivalence relation on
$U$ over $S$ (see
Groupoids, Definition \ref{groupoids-definition-equivalence-relation}).
\item the morphisms $s, t : R \to U$ are \'etale, and
\item the diagram
$$
\xymatrix{
R \ar@<1ex>[r] \ar@<-1ex>[r] &
U \ar[r] &
F
}
$$
is a coequalizer diagram in $\Sh((\Sch/S)_{fppf})$.
\end{enumerate}
\end{lemma}

\begin{proof}
Let $T/S$ be an object of $(\Sch/S)_{fppf}$.
Then $R(T) = \{(a, b) \in U(T) \times U(T) \mid f \circ a = f \circ b\}$
which defines an equivalence relation on $U(T)$.
The morphisms $s, t : R \to U$ are \'etale because the morphism
$U \to F$ is \'etale.

\medskip\noindent
To prove (3) we first show that
$U \to F$ is a surjection of sheaves, see
Sites, Definition \ref{sites-definition-sheaves-injective-surjective}.
Let $\xi \in F(T)$ with $T$ as above. Let $V = T \times_{\xi, F, f}U$.
By assumption $V$ is a scheme and $V \to T$ is surjective \'etale.
Hence $\{V \to T\}$ is a covering for the fppf topology.
Since $\xi|_V$ factors through $U$ by construction we
conclude $U \to F$ is surjective. Surjectivity implies that
$F$ is the coequalizer of the diagram by
Sites, Lemma \ref{sites-lemma-coequalizer-surjection}.
\end{proof}

\end{document}
